\documentclass[11pt]{article}

\usepackage[margin=1in]{geometry}
\usepackage{amsmath,amssymb,amsthm,mathtools}
\usepackage{booktabs,tabularx,array}
\usepackage{microtype}
\usepackage{enumitem}
\usepackage{xcolor}
\usepackage{xurl}
\usepackage[hidelinks]{hyperref}

% PORTFOLIO-RATIFICATION-PENDING: drafted 2026-07-12 (24h publication run)
% as the Paper E companion once the dynamics gate closed: generator layer
% (PlueckerPairGenerator), exact circuit locality (PlueckerCausalCone,
% PlueckerLayerCone), operational discriminator (PlueckerPhaseObservable),
% and the exact interacting two-particle spectrum fixture (pending
% harvest slots marked PENDING below). RELEASE GATES shared with the
% companions: named authors, artifact manifest.

\newcommand{\ii}{\mathrm{i}}
\newcommand{\id}{\mathbf{1}}
\newcommand{\C}{\mathbb{C}}
\newcommand{\tr}{\operatorname{tr}}
\definecolor{kernelcolor}{RGB}{0,92,74}
\definecolor{draftcolor}{RGB}{145,76,0}
\newcommand{\Kernel}{\textcolor{kernelcolor}{\textsc{Kernel}}}
\newcommand{\DraftTrust}{\textcolor{draftcolor}{\textsc{Kernel+Eval}}}

\newtheorem{theorem}{Theorem}
\newtheorem{remark}{Remark}

\title{Exact Pair-Gate Dynamics on a Fermionic Walk:\\
generator, circuit cones, and an exact interacting spectrum}
\author{Mark Schwab}
\date{July 12, 2026 --- working draft}

\begin{document}
\maketitle

\begin{abstract}
We assemble a largely kernel-checked account of the interacting layer of
a finite fermionic quantum walk, with the trust mark stated at each
result --- kernel where the statement is proved directly, and compiled
evaluator (on a Gaussian-integer twin) for the explicit $28\times28$
matrix identities behind the exact interacting spectrum.  The supplied two-particle phase kick ---
previously an opaque unitary gate --- is structured: it equals, up to a
global phase, the exact quarter-period pulse of a finite Hermitian even
quartic CAR generator $K(z)$ whose cube closes, $K^3=|z|^2K$, so the
generated one-parameter gate family closes into an explicit
circle-parametrized group, with the group law proved as pure algebra and an explicit
counterexample showing the naive half-pulse identification (without the
global phase) is false.  Locality is exact rather than asymptotic:
gates on disjoint mode sets commute, a scheduled sequence of finite-range
blocks keeps CAR support inside the corresponding iterated neighborhood,
and pairwise-disjoint layers cost one neighborhood step per layer, so
depth --- not gate count --- bounds propagation.  The free and kick
layers do not commute (four-mode witness, oracle-exact), so the composed
automaton is defined, as in the interacting
cellular-automaton literature, by alternating the free minor-lift with
the kick layer.  On an explicit four-site ring at the Pythagorean kick,
the composed step's two-particle characteristic polynomial is pinned
exactly (Section~4): the polynomial factorization into free levels and a
palindromic degree-12 factor, its reduction to a rational cubic in
$w=2\cos2\varepsilon$, and the identification with the actual composed
step are kernel; the underlying $28\times28$ Gaussian-integer matrix
arithmetic runs a disclosed compiled evaluator.  A
convention-fixed operational discriminator --- return probability $4/5$
versus $1$ at the displayed fixtures
$z=3+4\ii$ and $z=5$ --- separates the two equal-modulus conjugate
fields, consistent with the interacting sector retaining the Pluecker
phase.  Boundaries are stated exactly: the
generator is supplied, not derived from the free carrier (an open
problem shared with the Thirring-automaton literature); single-kick
spectra are phase-independent; and no thermodynamic or continuum claim
is made.
\end{abstract}

\section{Setup}

All objects are finite and explicit.  The one-particle space is
$\C^{2L}$ for a ring or line of $L$ sites with two internal components
per site; Fock space is the $2^{2L}$-dimensional occupation-basis space
over these modes, with creation and annihilation operators carrying the
standard ordered sign conventions (\texttt{FiniteCARFockBasic}: the CAR
relations, self-adjointness of the occupation form, and the sign
bookkeeping are kernel theorems, not assumptions).  A one-particle
unitary $U$ enters the many-body theory through the determinant-minor
lift $\Gamma(U)$, whose entries on the $n$-particle sector are the
$n\times n$ minors of $U$: the lift fixes the vacuum, preserves particle
number and parity, obeys exact creation covariance, and is functorial,
$\Gamma(1)=\id$ and $\Gamma(UV)=\Gamma(U)\Gamma(V)$, with the lifted
conjugate transpose an exact two-sided inverse
(\texttt{FiniteCARSecondQuantization}, \Kernel{}).  The supplied
interaction is an even quartic CAR polynomial supported on one pair
block (\texttt{PlueckerQuarticInteraction}, \Kernel{}); its structure is
the subject of Section~2.  Throughout, ``exact'' means an identity of
matrices over $\mathbb{Q}(\ii)$ checked by the Lean kernel; trust marks
follow the companion papers (\Kernel{} vs \DraftTrust{}).

\section{The generator layer}

\begin{theorem}[The pair-gate family closes into a circle group]\label{thm:gen}
Let $K(z)\psi = z\,(\text{pairForward}\,\psi) +
\bar z\,(\text{pairBackward}\,\psi)$ on the four-mode Fock space.  Then
$K$ is Hermitian for the occupation inner product; $K^3=|z|^2K$; the
circle-parametrized family
$U(c,s)=\id+(c-1)P-\ii s\,K/|z|$ (with $P$ the pair-block projector)
satisfies the exact group law
$U(c_1,s_1)U(c_2,s_2)=U(c_1c_2-s_1s_2,\,s_1c_2+c_1s_2)$ and is unitary
whenever $c^2+s^2=1$; and the walk's phase kick with unit coefficient
$z/|z|$ equals $\ii\,U(0,1)$ on the pair sector, both fixing the
complement --- the kick is the quarter-period pulse up to a global
phase.  The naive identification without that phase is false, by an
explicit amplitude witness.
\end{theorem}
\noindent\Kernel{}\quad(\texttt{PlueckerPairGenerator}:
\texttt{generator\_hermitian}, \texttt{generator\_cubed},
\texttt{group\_law}, \texttt{unitary}, \texttt{halfpulse\_low/high/off},
\texttt{naive\_halfpulse\_false}, all-rational witness and the
$(c,s)=(1,1)$ unitarity-failure control; in-file axiom guards.)

\section{Exact circuit locality}

Locality in this paper is a set of finite identities, not an asymptotic
bound.  Even gates on disjoint mode sets commute exactly, so a layer of
pairwise-disjoint gates is order-independent.  For a scheduled sequence
of finite-range blocks, CAR support propagates inside the corresponding
iterated graph neighborhood: the image of an operator supported on a
set $S$ under the scheduled circuit is supported on the union of the
blocks' neighborhoods reached by the schedule, with an explicit
contiguous witness realizing the bound and a nonlocal control showing
the hypothesis is load-bearing (\texttt{PlueckerCausalCone}, \Kernel{}).
A further layer theorem sharpens the count: pairwise-disjoint local
gates cost one neighborhood expansion per \emph{layer}, so an arbitrary
finite schedule is bounded by its layer depth rather than its gate
count (\texttt{PlueckerLayerCone}, \Kernel{}).  These are the
discrete-time, finite-volume counterparts of Lieb--Robinson-type
statements, but they are exact: no exponential tail, no velocity
constant to estimate.

\begin{remark}[No joint exponential]
$[H_{\mathrm{free}},K(z)]\neq0$ exactly (norm $5$ at the four-mode line
witness; oracle-exact, not formalized), so the free and kick layers do
not exponentiate jointly; we \emph{define} the interacting automaton by
alternation, matching the field convention, and within-layer exactness
is the landed theorem.  The ring-level witness is checked up to the same
compiled evaluator on the twin (\DraftTrust):
the pair-lifted translation and the composed step have an explicit
nonzero commutator entry
(\texttt{PairMomentumBlocks.kick\_breaks\_translation}, \DraftTrust{}).
\end{remark}

\section{The exact interacting spectrum on a four-site ring}
At the rational $3$-$4$-$5$ kick the composed step's two-particle
characteristic polynomial is pinned exactly.  A kernel polynomial
factorization (\texttt{PairSpectrumFixture.charpoly\_factorization},
over any commutative ring) gives the product
$(\lambda+1)^2(\lambda-1)^4(25\lambda^2+14\lambda+25)
(5\lambda^2-6\lambda+5)^2(5\lambda^2+6\lambda+5)^2\,p_{12}(\lambda)$
as the explicit degree-$28$ polynomial with leading coefficient $5^{11}$;
a kernel palindromic reduction (\texttt{p12\_palindromic\_reduction},
over any field) rewrites the interacting factor as
$p_{12}(\lambda)=\lambda^6\,(3125w^3-2300w^2-6156w-1440)$,
$w=\lambda^2+\lambda^{-2}$, so the twelve interaction-shifted
quasienergies solve the rational cubic in $w=2\cos2\varepsilon$; and the
matrix bridge (\texttt{PairCharpolyBridge.V\_charpoly\_eq}) proves
$5^{11}\,\chi_V(\lambda)$ equals that same explicit polynomial --- the
characteristic polynomial of the composed step \emph{is} the displayed
product, with the polynomial charpoly obtained structurally through a
kernel companion-matrix assembly rather than a symbolic determinant.
Six exact pinned modes (\texttt{Vz\_eigenvector\_plus/minus}, kernel
\texttt{decide}) --- four $+1$ and two $-1$ eigenvectors --- realize the
$(\lambda-1)^4$ and $(\lambda+1)^2$ free factors, and the faithfulness
identity $B_zK_z=5\,V_z$ with $V=U_2K_2$ ties the integer twin to the
physical matrix.  The polynomial-identity, structural-charpoly, and
pinned-mode steps are kernel; the twin-matrix arithmetic identities run
\texttt{native\_decide} on the $28\times28$ Gaussian-integer twin
(compiled-evaluator trust), verified in the Aristotle build environment
and pinned in a dedicated high-memory guard
(\texttt{PairSpectrumFixtureGuard}) rather than the always-on aggregate
guard, since the twin computation exceeds a $34$\,GB rebuild.  The
structural companion is checked up to a compiled evaluator on a
Gaussian-rational surrogate
(\texttt{PairMomentumBlocks}, \DraftTrust{}): its kernel statements are
transported from two disclosed compiled-evaluator identities on a
Gaussian-rational twin of the momentum-block operator, so this layer is
eval-trusted and about the twin rather than a kernel proof of the
actual-field operator.  The cyclic momentum $K\in\mathbb{Z}/4$ splits
the 28
two-particle modes into blocks of dimension $(6,8,6,8)$; the $K=0$ block
carries four $+1$ and two $-1$ modes, the $K=2$ block carries $2/2/2$
($+1$/$-1$/the doubled-phase pair $((3+4\ii)/5)^2$), and the $K=1,3$
blocks carry the two Pythagorean quadruples; so the composed automaton
admits no exact momentum decomposition: the kick has nonzero matrix
elements between distinct momentum sectors (explicit entry $3\ii/5$) and
weights the four sectors equally on the seed state
($|P_K\,e_{01}|^2=1/4$).  Whether every interacting quasienergy draws
support from every sector is an eigenvector-participation question we do
not resolve here.

\section{Operational phase retention}

The interacting sector can \emph{measure} the phase the free spectrum
cannot see.  Equal-modulus conjugate rest operators
produce return probabilities $4/5$ versus $1$ at the displayed fixtures
($z=3+4\ii$ versus $z=5$) in the supplied pair-kick
sector (\texttt{PlueckerPhaseObservable}, \Kernel{}); single-kick
spectra are exactly phase-independent, so the discriminator is genuinely
interferometric --- it requires two kicks and the relative phase between
them, exactly as a two-arm interferometer requires two beam splitters.

\section{Boundaries and open problems}
The generator is supplied, not derived --- open here and in the
interacting-automaton literature; the one-particle free carrier does
carry the transported relative phase spectrally (companion paper).
The companion paper pre-registers the sharp version as a conjecture
with gate and kill condition: among one-pair-block Hermitian even
quartic CAR polynomials, the phase-covariance constraint should select
exactly the real two-dimensional family of the displayed generator and
its quarter-phase rotation; the gate is finite rational linear algebra,
and a larger solution space kills the uniqueness reading while
promoting the surplus generators to results.
No continuum, thermodynamic, or scattering-theory claim is made; the
spectrum fixture is $L=4$.

\appendix
\section{Machine verification}
[PENDING: module/guard table on freeze.]

\begin{thebibliography}{9}
\bibitem{ThirringQCA}
A.~Bisio, G.~M.~D'Ariano, P.~Perinotti, and A.~Tosini,
\emph{The Thirring quantum cellular automaton},
Phys.\ Rev.\ A 97 (2018) 032132, arXiv:1711.03920.
\bibitem{ThirringPerturbative}
A.~Bisio, P.~Perinotti, A.~Pizzamiglio, and S.~Rota,
\emph{A perturbative approach to the solution of the Thirring quantum
cellular automaton}, Entropy 27 (2025) 198, arXiv:2406.19917.
\bibitem{QCAQFT}
L.~Mlodinow and T.~A.~Brun,
\emph{Quantum field theory from a quantum cellular automaton in one
spatial dimension and a no-go theorem in higher dimensions},
Phys.\ Rev.\ A 102 (2020) 042211, arXiv:2006.08927.
\end{thebibliography}

\end{document}
